\section{Methodology}
\subsection{How Efficient are Privacy Metrics?}
MIAs based on shadow modeling (smMIAs) \citep{stadler2022synthetic, houssiau2022tapas, meeus2023achilles} are the state-of-the-art and hence the most sensitive metric to privacy leaks in generative models. However, they are prohibitively expensive: they require training hundreds of generative models on different datasets, along with a large auxiliary dataset that often exceeds the private training data, making them impractical as privacy metrics in real-world scenarios. Furthermore, these attacks target a single record at a time, making batch evaluation even costlier. For instance, \citet{meeus2023achilles} uses datasets of size 1,000 but requires 75,000 records in total to evaluate the attack.
DOMIAS \citep{van2023membership} is more practical, requiring only a single synthetic dataset, but still needs an auxiliary dataset and involves training two density models, which can be expensive unless the density model is very simple (for instance, a kernel density estimator). For instance, in one of their experiments training and control sets each contain 500 records, and the auxiliary set has 10,000 records, implying the training set is 22 times smaller than the total amount of data used.
DCR-based metrics are instead more efficient, requiring no additional SDG calls and only a control dataset of the same size as the training data.
Despite the recent advances in MIA research, there remains a need for a practical privacy metric that avoids expensive model training and does not require auxiliary data far exceeding the training set.

\subsection{ReMIA}
We aim to design a practical MIA-based metric that does not suffer from the aforementioned limitations.
To do this, we draw inspiration from the idea of MIAs and we propose a new type of attack: we evaluate the privacy of a generative method by testing whether an attacker can determine if a specific record belongs to one of two distinct training datasets from which two synthetic datasets were generated by the generative model.
If the attacker can successfully determine the origin of the record with high accuracy, it indicates that the generator is not privacy-preserving, as the synthetic data leaks specific information about individual records in the training data.
This is in contrast with MIAs, which test whether an attacker can determine if a specific record was part of the training data or not.
This approach allows us to evaluate the privacy of the generative method without requiring an auxiliary dataset, and without training the generative model more than twice.

We call our method ReMIA, which stands for "Relative Membership Inference Attack", since it is based on the idea of identifying which training dataset a record belongs to, rather than guessing whether a record was part of the training data or not as in standard MIAs.
Similar to MIAs, we can formalize an attack as a privacy game between a challenger $C$, having access to a private dataset $\data$ and a generative model $\Phi$, and an attacker $A$:
\begin{enumerate}
    \item $C$ builds two training datasets $\train_1 = (\target_1, \shared)$ and $\train_2 = (\target_2, \shared)$, where $\target_1$, $\target_2$, and $\shared$ are disjoint sets of records sampled from the same private dataset $\data$, with $|\target_1| = |\target_2|$.
    \item $C$ trains $\Phi$ on $\train_1$ and $\train_2$ to generate two synthetic datasets $\synthetic_1 = \Phi(\train_1)$ and $\synthetic_2 = \Phi(\train_2)$.
    \item $C$ sends $\synthetic_1$, $\synthetic_2$ and $\target = (\target_1, \target_2)$ (shuffled, without label information) to the attacker $A$.
    \item $A$ uses $\synthetic_1$ and $\synthetic_2$ to build a classifier $\classifier$ that predicts whether each record in $\target$ belongs to $\target_1$ or $\target_2$. The score of the attack is the performance of the classifier $\classifier$ on $\target$.
\end{enumerate}
If the attacker can successfully determine the origin of the record with better-than-random accuracy, it indicates that the generator is not privacy-preserving, as the synthetic data leaks specific information about individual records in the training data.
On the other hand, if the attacker cannot distinguish between $\synthetic_1$ and $\synthetic_2$ with better than random guessing, it suggests that the generator is effectively preserving privacy by generating synthetic data that does not depend on specific records in the training data.
Throughout the paper, we use the AUROC of the classifier $\classifier$ as the privacy score of ReMIA.

\subsection{An Attacker Model for ReMIA}
Building an attacker model for ReMIA is not straightforward, as we do not assume the attacker to have access to auxiliary data or the generative model.
We propose a simple method consisting in training a classifier $\classifier$ to distinguish between the two synthetic datasets $\synthetic_1$ and $\synthetic_2$, and then evaluating the performance of this classifier on the two unseen sets of records $\target_1$ and $\target_2$.

In practice, in order to observe a signal on unseen data $\target_1$ and $\target_2$, we train the classifier $\classifier$ past the point of overfitting on the synthetic data $\synthetic_1$ and $\synthetic_2$, and we stop when a certain threshold of performance on the training synthetic data is reached. For instance, we stop when the AUROC reaches 99\%.
Notice that this is not a problem, since we are not interested in the performance of the classifier on the synthetic data, but rather in its performance on the unseen data $\target_1$ and $\target_2$ as a "side effect".
The intuition behind this attack method is that if the classifier can successfully distinguish between $\synthetic_1$ and $\synthetic_2$, it may be able to learn features that are specific to the training datasets $\train_1$ and $\train_2$, which in turn may allow it to correctly classify records in $\target_1$ and $\target_2$ based on the synthetic data.
In order to reduce the noise of the privacy score, we apply a smoothing function to the performance of the classifier on $\target_1$ and $\target_2$ across the learning iterations, and select as the score the smoothed value at the iteration where the performance on the synthetic data reaches a certain threshold.

Notice that without having access to the labels of $\target$, the attacker cannot know at which training iteration the classifier reaches the best performance on the unseen data $\target$, as we show in the example in Figure \ref{remia_score_learning_step} in Appendix \ref{app:additional_results}. The stopping criterion based on the performance on the synthetic data is a practical heuristic that we will show to be effective in practice. In principle, one can select the maximum performance on the unseen data $\target$ as the privacy score. However, this can lead to a biased estimate of the privacy risk. This can be solved by assuming the attacker to have access to a subset of labeled records from $\target$ that can be used as a validation set to select the best iteration. We leave this analysis for future work.

We implement the classifier $\classifier$ as a multi-layer perceptron (MLP). More details about the architecture, hyperparameters, smoothing function, and the training procedure are provided in the experimental section and Appendix \ref{app:remia_details}.



\subsection{Efficiency of ReMIA}
Notice that ReMIA only requires training the generative model twice (on $\train_1$ and $\train_2$), it does not require any auxiliary dataset, but only a hold-out split of the original dataset, and the attacker model only consists in training a single classifier.
The data usage of ReMIA is also efficient. Let $s = |\train_1| = |\train_2|$ be the size of each training set, where $t = |\target_1| = |\target_2|$ is the number of target records in each training set, and $r = |\shared|$ is the number of shared records so that $s = t + r$. If the fraction of target records in each training dataset is $f = \frac{t}{s}$, then the total number of records needed to run ReMIA is $n = |\target_1|+|\target_2|+|\shared| = s(1+f)$, which is at most twice the size of the training dataset used to train the generative model.
This allows us to evaluate the privacy of the generative model on training datasets whose size is comparable to the original dataset. 
We can choose a fraction of target records $f$ that is large enough to obtain a statistically stable result, i.e., that does not depend on the specific set of target records. On the other hand, a small fraction allows us to evaluate ReMIA on a training dataset that is closer in size to the total amount of available data.

While smMIAs represent powerful but prohibitive privacy metrics, ReMIA achieves practicality by trading off the attack realism for computational efficiency.
We summarize the ReMIA privacy score computation in Algorithm \ref{alg:remia}.
\begin{algorithm}[t]
\caption{ReMIA privacy score computation}
\label{alg:remia}
\begin{algorithmic}[1]
\Require Dataset $\data$, generative model $\Phi$, targets fraction $f$, iterations $N$
\Ensure ReMIA privacy score
\State Randomly Split $\data$ into $\target_1$, $\target_2$, and $\shared$, with $|\target_1| = |\target_2| = \frac{f}{1+f}|\data|$ and $|\shared| = |\data| - 2|\target_1|$
\State Create two training datasets $\train_1 = \target_1 \cup \shared$ and $\train_2 = \target_2 \cup \shared$
\State Generate synthetic datasets $\synthetic_1 = \Phi(\train_1)$ and $\synthetic_2 = \Phi(\train_2)$
\State Create a dataset $\synthetic = \synthetic_1 \cup \synthetic_2$ and $\target = \target_1 \cup \target_2$ with origin labels for each record
\State Initialize a classifier $\classifier$
\For{$i = 1$ \textbf{to} $N$}
  \State Train $\classifier$ on $\synthetic$ for one iteration
  \State Evaluate AUROC of $\classifier$ on $\target$ to obtain $P_{\text{target}}^i$ and on $\synthetic$ to obtain $P_{\text{train}}^i$
\EndFor
\State Apply a smoothing function to scores $\{P^i_{\text{target}}\}_{i=1}^N$ to obtain smoothed scores $\{\tilde{P}^i_{\text{target}}\}_{i=1}^N$
\State \Return $\tilde P^k_{\text{target}}$ where $k = \min\{ i \mid P_{\text{train}}^i \geq 99\%\}$
\end{algorithmic}
\end{algorithm}
